\chapter{Complexity — IS Count Gen Machine}
%%% generated by the blueprint pipeline from TCSlib.Complexity.SpaceComplexity.Logspace.ISCountGenMachine
\section{Overview}

Declarations of the Lean module \texttt{TCSlib.Complexity.SpaceComplexity.Logspace.ISCountGenMachine}.

\section{Declarations}

\begin{definition}[Logarithmic counting control states]
\lean{Complexity.LogProg.ISGen.GLb}
\label{Complexity.LogProg.ISGen.GLb}
\leanok
For each natural number $E$, introduce a finite set of control states for a counting
procedure. Its states represent the setup and logarithm calculation, parameterized
rounds and doubling, counting and nested enumeration stages, acceptance or rejection,
and a walk through a sequence of configurations; states within a round carry a Boolean
phase and an index among $0,\ldots,E$. The walk states carry Boolean flags as indicated
by their roles, including whether to query an initial-code or step decider, compare
against a target or bound, and count a step.
\end{definition}

\begin{definition}[Walk comparison state]
\lean{Complexity.LogProg.ISGen.gwX}
\label{Complexity.LogProg.ISGen.gwX}
\leanok
\uses{Complexity.LogProg.ISGen.GLb}
For each natural number $E$ and Boolean value $b$, define a control state that is the
inner-enumeration equality-check state when $b$ is true, and the outer-enumeration
accepting-code query state when $b$ is false.
\end{definition}

\begin{definition}[Walk continuation state]
\lean{Complexity.LogProg.ISGen.gwcX}
\label{Complexity.LogProg.ISGen.gwcX}
\leanok
\uses{Complexity.LogProg.ISGen.GLb}
For each natural number $E$ and Boolean values $b,f$, define a control state that counts
a step when $f$ is true and otherwise clears the step counter; the state carries the
value $b$.
\end{definition}

\begin{definition}[Round initialization instruction]
\lean{Complexity.LogProg.ISGen.grdX}
\label{Complexity.LogProg.ISGen.grdX}
\leanok
\uses{Complexity.LogProg.ISCount.rK, Complexity.LogProg.ISCount.rk, Complexity.LogProg.ISGen.GLb, Complexity.LogProg.Ins}
For each natural number $E$ and Boolean phase $p_h$, define a register-machine
instruction on 17 registers and four deciders: if $p_h$ is true, clear register $K$ and
continue at the state that sets $K$ to one; otherwise, clear register $k$ and continue
at the first round state, with the phase set to true and the round index set to $E$.
\end{definition}

\begin{definition}[Counting procedure instruction map]
\lean{Complexity.LogProg.ISGen.igA}
\label{Complexity.LogProg.ISGen.igA}
\leanok
\uses{Complexity.LogProg.ISCount.rK, Complexity.LogProg.ISCount.rT, Complexity.LogProg.ISCount.rY, Complexity.LogProg.ISCount.rc, Complexity.LogProg.ISCount.rcnt, Complexity.LogProg.ISCount.rcp, Complexity.LogProg.ISCount.rcur, Complexity.LogProg.ISCount.re, Complexity.LogProg.ISCount.rg, Complexity.LogProg.ISCount.ri, Complexity.LogProg.ISCount.ri1, Complexity.LogProg.ISCount.rk, Complexity.LogProg.ISCount.rlg, Complexity.LogProg.ISCount.rn, Complexity.LogProg.ISCount.rs, Complexity.LogProg.ISCount.rstp, Complexity.LogProg.ISCount.ru, Complexity.LogProg.ISCount.rv, Complexity.LogProg.ISCount.wB, Complexity.LogProg.ISCount.wT, Complexity.LogProg.ISGen.GLb, Complexity.LogProg.ISGen.grdX, Complexity.LogProg.ISGen.gwX, Complexity.LogProg.ISGen.gwcX, Complexity.LogProg.NARM}
For each natural number $E$, define a nondeterministic register-machine instruction map
with 17 registers, four deciders, and the control states indexed by $E$. It initializes
the input-length and logarithm registers, computes its round parameters, and then
performs the bounded counting and enumeration procedure: it may choose candidates for
the enumerations and walk, queries the specified deciders, updates counters and
registers, and ultimately returns true or false according to the designated acceptance
and rejection states.
\end{definition}

\begin{lemma}[Instruction Safety of the Generated Program]
\lean{Complexity.LogProg.ISGen.igA_npreS}
\label{Complexity.LogProg.ISGen.igA_npreS}
\leanok
\uses{Complexity.LogProg.AConf, Complexity.LogProg.ISCount.rK, Complexity.LogProg.ISCount.rT, Complexity.LogProg.ISCount.rY, Complexity.LogProg.ISCount.rc, Complexity.LogProg.ISCount.rcnt, Complexity.LogProg.ISCount.rcp, Complexity.LogProg.ISCount.rcur, Complexity.LogProg.ISCount.re, Complexity.LogProg.ISCount.rg, Complexity.LogProg.ISCount.ri, Complexity.LogProg.ISCount.ri1, Complexity.LogProg.ISCount.rk, Complexity.LogProg.ISCount.rlg, Complexity.LogProg.ISCount.rn, Complexity.LogProg.ISCount.rs, Complexity.LogProg.ISCount.rstp, Complexity.LogProg.ISCount.ru, Complexity.LogProg.ISCount.rv, Complexity.LogProg.ISCount.wB, Complexity.LogProg.ISCount.wT, Complexity.LogProg.ISGen.GLb, Complexity.LogProg.ISGen.grdX, Complexity.LogProg.ISGen.igA, Complexity.LogProg.NARM.detPart, Complexity.LogProg.NPreS, Complexity.LogProg.PreS}
\difficulty{5}
For every natural number $E$, every finite Boolean word $x$, and every configuration of
the generated register program, the instruction safety conditions hold: an active
register-equality test compares distinct registers; a test comparing a register with the
input payload occurs only when $x$ is a valid plain encoding; tests comparing a register
with either inner payload occur only when $x$ is a valid paired encoding; and every
decider call has a list of pairwise distinct argument registers. A halted configuration
or a configuration whose active instruction is a nondeterministic choice satisfies the
condition without further restrictions.
\end{lemma}

\begin{definition}[Bounded configuration semantics]
\lean{Complexity.LogProg.ISGen.IG}
\label{Complexity.LogProg.ISGen.IG}
\leanok
\uses{Complexity.LogProg.ISCount.St, Complexity.LogProg.ISCount.St.toV, Complexity.LogProg.ISGen.GLb, Complexity.LogProg.ISGen.igA, Complexity.LogProg.NSem, Complexity.LogProg.RegLe}
For a natural number $E$, an oracle
$\mathrm{orc}:\mathrm{Fin}(4)\to\mathrm{List}(\mathrm{Bool})\to\mathrm{Bool}$, an input
$x\in\mathrm{List}(\mathrm{Bool})$, a natural-number bound $B$, a state label
$l\in\mathrm{GLb}(E)$, a register state $\sigma$, and a proposition $P$, the predicate
holds exactly when the nondeterministic semantics of the machine parameterized by $E$,
with oracle $\mathrm{orc}$ and register bound $\mathrm{RegLe}(17,\mathrm{GLb}(E),B)$,
assigns $P$ to the initial configuration with control state $l$, register vector
$\sigma$, and no output.
\end{definition}

\begin{lemma}[Register bound from state bounds]
\lean{Complexity.LogProg.ISGen.St.greg}
\label{Complexity.LogProg.ISGen.St.greg}
\leanok
\uses{Complexity.LogProg.ISCount.St, Complexity.LogProg.ISCount.St.Bd, Complexity.LogProg.ISCount.St.toV, Complexity.LogProg.ISGen.GLb, Complexity.LogProg.RegLe}
\difficulty{4}
Let $E$ and $B$ be natural numbers, let $\sigma$ be a counting-algorithm state whose
seventeen registers are each at most $B$, let $l$ be an optional control label, and let
$res$ be an optional Boolean result. The abstract configuration with label $l$, register
vector from $\sigma$, and result $res$ has every register at most $B$.
\end{lemma}

\begin{lemma}[Invariance under state equality]
\lean{Complexity.LogProg.ISGen.IG.of_eq}
\label{Complexity.LogProg.ISGen.IG.of_eq}
\leanok
\uses{Complexity.LogProg.ISCount.St, Complexity.LogProg.ISGen.GLb, Complexity.LogProg.ISGen.IG}
\difficulty{1}
Fix a Boolean oracle on four indices and finite Boolean strings, an input string, a
natural-number register bound $B$, a control label $l$, and a proposition $P$. If two
counting-algorithm states $\sigma$ and $\sigma'$ are equal, and the abstract machine
starting at $l$ with register vector from $\sigma'$ terminates with all registers at
most $B$ and accepts if and only if $P$ holds, then the same bounded termination and
acceptance property holds when it starts with the register vector from $\sigma$.
\end{lemma}
